% ===================================================================
%  Insights into Imaging — Main Text
%  Scanner Vendor Harmonisation of Spinal Cord qMRI Biomarkers
% ===================================================================
\documentclass[12pt,a4paper]{article}

% Packages
\usepackage[utf8]{inputenc}
\usepackage[T1]{fontenc}
\usepackage{amsmath,amssymb}
\usepackage{graphicx}
\usepackage{booktabs}
\usepackage[margin=2.5cm]{geometry}
\usepackage{hyperref}
\usepackage{setspace}
\usepackage{caption}
\usepackage{enumitem}
\usepackage{float}
\usepackage{lineno}
\usepackage{multirow}
\usepackage{siunitx}
\usepackage{tabularx}
\usepackage{pdflscape}

\hypersetup{colorlinks=true,linkcolor=blue,citecolor=blue,urlcolor=blue}
\onehalfspacing

\begin{document}
\linenumbers

% ===================================================================
\section*{Abstract}

\noindent\textbf{Objectives:} Quantitative spinal cord MRI biomarkers exhibit systematic vendor-dependent variability that limits multi-centre pooling, yet no systematic harmonisation benchmark exists for spinal cord data.

\noindent\textbf{Methods:} Eight qMRI biomarkers (T2w-CSA, GM-CSA, MTR, MTsat, FA, MD, AD, RD) were extracted at C2--C3 from the spine-generic dataset (267 subjects, 44 sites, 3 vendors; site and vendor perfectly nested). Five methods---ComBat, ComBat-joint, CovBat, RELIEF, LME/FE---were evaluated on three endpoints in healthy-control ($N=188$) and pathology-inclusive ($N=246$) cohorts: vendor-effect removal, biological-signal recovery, and subject-level preservation. Leave-one-vendor-out validation tested generalisability.

\noindent\textbf{Results:} Baseline vendor effects were substantial ($\eta^2=0.28$; PERMANOVA $R^2=0.32$, $p<10^{-4}$). All methods reduced vendor $\eta^2$ by $>99.8\%$ (to ${\sim}10^{-4}$), rendering multivariate vendor structure non-significant (all $p\geq0.97$). Age signal increased 2.3--3.7-fold. Methods diverged on subject preservation (LME/FE: $r=0.84$; CovBat: $r=0.78$), forming a Pareto frontier. LOVO validation showed ComBat/ComBat-joint generalise perfectly ($\eta^2<10^{-3}$, PERMANOVA $p>0.86$), while RELIEF and LME/FE failed ($\eta^2\sim0.19$, $p<0.03$). Three methodological constraints---vendor-site nesting, LME fixed-effects fallback (75\% of cases), and RELIEF ICA inactivation---converge on a single structural limitation: only 3 vendors in a perfectly nested 44-site design.

\noindent\textbf{Conclusion:} Harmonisation is recommended for pooled multi-vendor spinal cord qMRI. ComBat provides robust generalisability and is preferred for broad deployment; CovBat/RELIEF are suitable for pooled-group analyses within fixed vendor sets.

\vspace{4mm}
\noindent\textbf{Key Points:}
\begin{enumerate}[nosep,leftmargin=*]
  \item \textbf{Question:} Multi-centre spinal cord quantitative MRI biomarkers exhibit systematic vendor-dependent variability that limits cross-site pooling, yet no harmonisation benchmark exists to guide method selection.
  \item \textbf{Findings:} All five methods removed over 99.8\% of vendor effects; ComBat generalised best in leave-one-vendor-out validation while preserving biological signal and subject-level ranking.
  \item \textbf{Critical Relevance Statement:} This benchmark provides the first quantitative evidence for selecting harmonisation methods in multi-centre spinal cord MRI, enabling reliable biomarker pooling across scanner vendors for clinical research and trials.
\end{enumerate}

\vspace{4mm}
\noindent\textbf{Keywords:} Spinal cord; Quantitative MRI; Harmonisation; ComBat; Multicentre study; Reproducibility; Scanner vendor effects

\vspace{4mm}
\noindent\textbf{Abbreviations:} AD = axial diffusivity; ComBat = Combatting batch effects; CovBat = Covariance Batch harmonisation; CSA = cross-sectional area; DCM = degenerative cervical myelopathy; EB = empirical Bayes; FA = fractional anisotropy; HC = healthy control; ICA = independent component analysis; LME/FE = linear mixed-effects / fixed-effects hybrid; LOVO = leave-one-vendor-out; MD = mean diffusivity; MT = magnetisation transfer; MTR = magnetisation transfer ratio; MTsat = magnetisation transfer saturation; PERMANOVA = permutational multivariate analysis of variance; qMRI = quantitative magnetic resonance imaging; RD = radial diffusivity; RELIEF = Removal of Latent Inter-scanner Effects via independent components

% ===================================================================
\section{Introduction}

Quantitative MRI (qMRI) biomarkers of the spinal cord---cross-sectional area (CSA), magnetisation transfer ratio (MTR), magnetisation transfer saturation (MTsat), and diffusion tensor metrics (FA, MD, AD, RD)---provide non-invasive windows into atrophy, demyelination, and axonal integrity in multiple sclerosis, degenerative cervical myelopathy, and spinal cord injury [1--3]. Multi-centre consortia are increasingly adopting these biomarkers, but systematic vendor-dependent variability limits cross-site pooling [4,5]. A recent pan-European survey by the European Society of Radiology documented marked heterogeneity in MRI quality assurance and control practices across 37 countries [17], underscoring the urgency of robust harmonisation.

Even under standardised acquisition protocols, differences in gradient hardware, RF coil geometry, and vendor-specific reconstruction pipelines introduce offsets that often exceed the biological effect sizes of interest [4]. Statistical post-hoc harmonisation methods---ComBat [6,7,16], CovBat [8], RELIEF [9], and linear mixed-effects models (LME) [10]---have been extensively benchmarked in brain imaging and reviewed in the broader context of multicentre radiomics [18,21] but remain uncharacterised for spinal cord qMRI. Fortin et al.\ first applied ComBat to multi-site DTI [16], establishing a framework later extended to cortical thickness [7], longitudinal harmonisation [19], covariance harmonisation [8], latent-factor removal [9], and multicentre clinical radiomics applications [20]. Middleton et al.\ demonstrated ComBat's effectiveness for spinal cord DTI [11], and Jodoin et al.\ provided deployment guidance for diffusion MRI [12], but no study has systematically compared the full ComBat family across a comprehensive spinal cord biomarker panel using unified multi-endpoint evaluation.

The spinal cord presents distinct harmonisation challenges: small per-site sample sizes (5--15 subjects) challenge empirical Bayes shrinkage; strong inter-metric correlations (diffusion metrics are intrinsically coupled; MTR and MTsat share acquisition-dependent variance) motivate multivariate evaluation; and the clinical need to preserve disease-related variance creates a tension not previously quantified.

Here we benchmark five methods on eight spinal cord qMRI biomarkers using three orthogonal endpoints: vendor-effect removal, biological-signal recovery, and subject-level preservation. We extend prior work through: (i) formal comparison of the full ComBat family (including covariance and latent-factor extensions); (ii) multivariate evaluation via PERMANOVA; (iii) a Pareto-frontier decision framework; and (iv) leave-one-vendor-out external validation to test generalisability.

% ===================================================================
\section{Materials and Methods}

\subsection{Dataset and Cohorts}

The spine-generic multi-subject dataset (release r20231212) [1] comprises 267 subjects scanned at 44 sites across three vendors (GE Healthcare: 37 subjects at 6 sites; Philips: 50 subjects at 8 sites; Siemens Healthineers: 180 subjects at 30 sites) using a standardised consensus protocol [4]. Site and vendor are perfectly nested (each site operates a single scanner from one vendor); reported vendor effects thus represent the composite ``vendor-in-site'' effect. Demographics were balanced across vendors (Table~1; all $p\geq0.125$).

\begin{table}[H]
\centring
\caption{Demographic and acquisition characteristics by vendor (HC cohort, $N=203$; complete-case $N=188$). Continuous variables: mean $\pm$ SD, one-way ANOVA. Categorical variables: count (\%), $\chi^{2}$ test. Site and vendor are perfectly nested (44 sites across 3 vendors).}
\label{tab:demographics}
\begin{tabular}{lrrrrr}
\toprule
Characteristic & GE & Philips & Siemens & Total & $p$-value \\
\midrule
$N$ & 28 & 36 & 139 & 203 & --- \\
Sites ($n$) & 6 & 8 & 30 & 44 & --- \\
Age, years (mean $\pm$ SD) & 26.9 $\pm$ 4.7 & 28.2 $\pm$ 4.5 & 29.1 $\pm$ 5.9 & 28.7 $\pm$ 5.6 & 0.125 \\
Age range & 19--39 & 23--38 & 20--52 & 19--52 & --- \\
Sex = F & 15 (53.6\%) & 19 (52.8\%) & 64 (46.0\%) & 98 (48.3\%) & 0.643 \\
Sex = M & 13 (46.4\%) & 17 (47.2\%) & 75 (54.0\%) & 105 (51.7\%) & --- \\
Height, cm & 169.6 $\pm$ 8.5 & 172.2 $\pm$ 10.2 & 173.2 $\pm$ 10.0 & 172.5 $\pm$ 9.9 & 0.232 \\
Weight, kg & 65.2 $\pm$ 13.7 & 67.6 $\pm$ 14.3 & 68.2 $\pm$ 13.1 & 67.7 $\pm$ 13.4 & 0.576 \\
Scanner models ($n$) & 3 & 4 & 6 & 13 & --- \\
\bottomrule
\end{tabular}
\end{table}

Two pre-specified cohorts were analysed: HC ($N=188$ complete-case healthy controls) isolates pure vendor-driven variance; ALL ($N=246$ complete-case, including 56 mild cervical compression and 2 DCM) approximates real-world deployment. Pathology distribution was balanced across vendors ($\chi^2\ p=0.661$). Detailed cohort flow and missingness patterns are provided in Supplementary Material~S1.

\subsection{Biomarker Extraction and Processing}

Eight biomarkers spanning four acquisition families were extracted at the C2--C3 level using the Spinal Cord Toolbox (SCT v7.4) [13]: T2w-CSA and GM-CSA (T2-weighted structural), MTR and MTsat (magnetisation transfer), and FA, MD, AD, RD (diffusion tensor imaging). Processing followed the spine-generic pipeline documentation. Automated quality-control metrics excluded acquisitions with segmentation or model-fit failures (see Supplementary Material~S2).

\subsection{Harmonisation Methods}

Five methods were evaluated (full mathematical details in Supplementary Material~S3):

\textbf{ComBat} [6,7,16]: Per-metric location-and-scale adjustment with empirical Bayes shrinkage on site-specific parameters. Each biomarker harmonised independently.

\textbf{ComBat-joint} [7]: All eight biomarkers harmonised jointly (joint empirical Bayes), preserving cross-metric covariance structure.

\textbf{CovBat} [8]: ComBat-joint followed by principal component (PC) correction---residuals projected onto PCs, each PC harmonised independently, then projected back.

\textbf{RELIEF} [9]: ComBat-joint followed by independent component analysis (ICA) of residuals. Components with significant vendor association (Benjamini-Hochberg FDR-controlled ANOVA, $q<0.05$) are zeroed. In the present dataset, no component reached significance (all $q\geq0.99$), rendering RELIEF functionally equivalent to ComBat-joint with additional PCA/ICA round-trip noise.

\textbf{LME/FE}: Site as a random intercept with age and sex fixed effects. When the random-effects covariance matrix was singular (expected with only 3 vendor levels---well below the recommended $k\geq5$ for stable REML estimation), the model fell back to ordinary least squares centring (75\% of 16 biomarker$\times$cohort combinations; see Results). Only MD and AD achieved stable LME fits in both cohorts. We label this hybrid approach ``LME/FE'' to acknowledge the dual mechanism.

\subsection{Evaluation Endpoints}

Three orthogonal endpoints were evaluated per method per cohort (detailed in Supplementary Material~S4):

\textbf{Vendor-effect removal}: Univariate $\eta^2$ (ANOVA), intraclass correlation coefficient (ICC), and multivariate PERMANOVA $R^2$ (Euclidean distance, 4,999 permutations) [14,15].

\textbf{Biological-signal recovery}: Type-II semi-partial $R^2$ for age and sex from linear models.

\textbf{Subject-level preservation}: Pearson $r$ between pre- and post-harmonisation per-subject biomarker vectors.

\subsection{Leave-One-Vendor-Out External Validation}

For each method, harmonisation models were trained on subjects from two vendors and applied to the held-out third vendor in three rotation folds. Test-vendor transformation used z-score distribution matching (ComBat, ComBat-joint) or projection onto training PCA/ICA space (CovBat, RELIEF). LME/FE assigned offset=0 to unseen vendors. The same three endpoints were evaluated on test-vendor data only (Supplementary Material~S5).

\subsection{Supplementary Analyses}

Design C balanced-subsample analysis drew 10 subjects per vendor (HC cohort, 1,000 bootstrap iterations) to isolate the effect of vendor imbalance from method performance. Full results in Supplementary Material~S6.

\subsection{Reporting Guidelines and Missing Data}

This study follows the STROBE (Strengthening the Reporting of Observational Studies in Epidemiology) guideline for cross-sectional studies where applicable. The spine-generic dataset is a curated benchmark with pre-processed, quality-controlled biomarker measurements. Subjects with missing biomarker data were excluded via complete-case analysis (HC: $N=188/203$, 7.4\%; ALL: $N=246/267$, 7.9\%); missingness was attributable to acquisition-level exclusions (motion artefacts, segmentation failures, and incomplete sequences) as detailed in Supplementary Material~S1. Sensitivity to missing-data assumptions was assessed through the ALL-cohort analysis, which preserves a broader subject spectrum.

% ===================================================================
\section{Results}

\subsection{Baseline Vendor Effects}

At baseline, vendor explained 28\% of univariate biomarker variance (mean $\eta^2=0.28$, range 0.02--0.71 across metrics) and 32\% of multivariate structure (PERMANOVA $R^2=0.32$, $p<10^{-4}$, HC cohort). These effects exceeded age-related variance by approximately two orders of magnitude and sex-related variance by one order. GM-CSA showed the largest vendor effect ($\eta^2_{\text{HC}}=0.71$, Philips scanners systematically lower; Figure~1a). In the ALL cohort, baseline vendor effects were comparable ($\eta^2=0.27$; PERMANOVA $R^2=0.30$).

\begin{figure}[H]
\centring
\includegraphics[width=\textwidth]{figures/Fig1.pdf}
\caption{Baseline vendor effects in spinal cord quantitative MRI (qMRI) biomarkers (healthy control cohort, $N=188$). \textbf{(a)} Box plots showing distributions of eight biomarkers---T2w cross-sectional area (T2w-CSA), grey matter CSA (GM-CSA), magnetisation transfer ratio (MTR), magnetisation transfer saturation (MTsat), fractional anisotropy (FA), mean diffusivity (MD), axial diffusivity (AD), radial diffusivity (RD)---by scanner vendor (GE, Philips, Siemens). \textbf{(b)} Principal coordinate ordination from permutational multivariate analysis of variance (PERMANOVA) using Euclidean distance, with points coloured by vendor.}
\label{fig:baseline}
\end{figure}

\subsection{Vendor-Effect Removal}

All five methods reduced vendor $\eta^2$ by approximately three orders of magnitude, to ${\sim}10^{-4}$ (median post-harmonisation $\eta^2$: ComBat $=8.9\times10^{-5}$, ComBat-joint $=4.5\times10^{-4}$, CovBat $=9.2\times10^{-5}$, RELIEF $=9.0\times10^{-5}$, LME/FE $=1.8\times10^{-6}$; Figure~2a--c, Table~2). Multivariate vendor structure became indistinguishable from chance (all PERMANOVA $p\geq0.97$ in HC cohort; $p\geq0.97$ in ALL cohort).

\begin{figure}[H]
\centring
\includegraphics[width=\textwidth]{figures/Fig2.pdf}
\caption{Harmonisation performance by method and endpoint (healthy control cohort). \textbf{(a)} Vendor effect size ($\eta^2$, eta-squared from one-way ANOVA) before and after harmonisation, shown on logarithmic scale. \textbf{(b)} Multivariate vendor structure (PERMANOVA $R^2$, permutational multivariate analysis of variance) with significance indicators (asterisks denote *$p<0.05$, **$p<0.01$, ***$p<0.001$; ns = not significant). \textbf{(c)} Age-related biological signal (semi-partial $R^2$) before and after harmonisation.}
\label{fig:performance}
\end{figure}

\begin{table}[H]
\centring
\caption{Harmonisation performance summary (HC cohort, $N=188$). Values are means across 8 biomarkers. $\eta^2_{\text{vendor}}$: post-harmonisation vendor ANOVA effect size; $R^2_{\text{PERMANOVA}}$: multivariate vendor structure; $r_{\text{pre,post}}$: within-vendor subject-ranking preservation; Age $R^2$: semi-partial age variance explained. LOVO columns show leave-one-vendor-out external validation (mean across 3 held-out-vendor folds).}
\label{tab:results}
\begin{tabular}{lSSSSSS}
\toprule
Method & {$\eta^2_{\text{vendor}}$} & {$R^2_{\text{PERMANOVA}}$} & {$r_{\text{pre,post}}$} & {Age $R^2$} & {LOVO $\eta^2$} & {LOVO ratio} \\
\midrule
Original (baseline) & 2.79e-01 & 0.315 & 1.000 & 0.003 & {---} & {---} \\
ComBat & 3.75e-04 & 0.002 & 0.817 & 0.006 & 1.72e-04 & 0.5$\times$ \\
ComBat-joint & 4.91e-04 & 0.001 & 0.806 & 0.009 & 4.42e-04 & 0.9$\times$ \\
CovBat & 2.54e-04 & 0.0003 & 0.782 & 0.010 & 1.01e-03 & 4.0$\times$ \\
RELIEF & 4.36e-04 & 0.0004 & 0.791 & 0.009 & 1.87e-01 & 428$\times$ \\
LME/FE & 3.07e-04 & 0.002 & 0.840 & 0.007 & 1.74e-01 & 565$\times$ \\
\bottomrule
\end{tabular}
\end{table}

\subsection{Biological-Signal Recovery}

Age $R^2$ increased 2.3--3.7-fold after harmonisation (from $8.7\times10^{-4}$ to median $1.9\times10^{-3}$--$3.2\times10^{-3}$ in HC; from $6.4\times10^{-3}$ to $1.5\times10^{-2}$--$2.4\times10^{-2}$ in ALL). Sex $R^2$ showed a similar pattern (1.3--2.2-fold increase). The amplification was larger in ALL, consistent with greater biological heterogeneity. Absolute age effects remained below 1\% of variance.

\subsection{Subject-Level Preservation and the Pareto Frontier}

Methods diverged on within-vendor subject ranking preservation (HC cohort, mean $r$ across methods: LME/FE $=0.840$, ComBat $=0.829$, ComBat-joint $=0.821$, RELIEF $=0.803$, CovBat $=0.780$). This ranking was stable between HC and ALL cohorts (ALL: 0.830, 0.826, 0.816, 0.793, 0.771, respectively). A clear Pareto frontier emerged: as multivariate decontamination (PERMANOVA $R^2$) increases, subject-level preservation decreases monotonically (Figure~3). Under balanced conditions (Design C), preservation differences compressed but maintained the same rank order (CovBat: $r=0.991$, others $r\geq0.998$; LME/FE: $r=0.841$), confirming that the preservation gradient reflects real-world vendor imbalance, not intrinsic method limitations.

\begin{figure}[H]
\centring
\includegraphics[width=\textwidth]{figures/Fig3.pdf}
\caption{Pareto frontier showing the trade-off between multivariate vendor-effect removal (PERMANOVA $R^2$, permutational multivariate analysis of variance) and subject-level preservation (Pearson correlation $r_{\text{pre,post}}$ between pre- and post-harmonisation biomarker vectors). Healthy control cohort. Each point represents one harmonisation method. Lower right indicates maximal vendor-effect removal; upper left indicates maximal subject-level preservation.}
\label{fig:pareto}
\end{figure}

\subsection{LOVO External Validation}

ComBat and ComBat-joint generalised to unseen manufacturers with negligible degradation (LOVO $\eta^2$: $1.7\times10^{-4}$ and $4.4\times10^{-4}$; PERMANOVA $p=0.864$ and $0.951$; LOVO/internal $\eta^2$ ratio: 0.5 and 0.9). CovBat showed moderate degradation ($\eta^2=1.0\times10^{-3}$, ratio 4.0). RELIEF and LME/FE failed external validation ($\eta^2_{\text{LOVO}}$: 0.19 and 0.17; PERMANOVA $p<0.03$; ratios 428 and 565; Figure~4a--d). This pattern is explained by structural constraints: RELIEF's ICA unmixing matrix is non-transferable across vendor sets, and LME/FE cannot estimate offsets for unseen vendors.

\begin{figure}[H]
\centring
\includegraphics[width=\textwidth]{figures/Fig4.pdf}
\caption{Leave-one-vendor-out (LOVO) external validation. \textbf{(a)} Vendor effect size ($\eta^2$, eta-squared from one-way ANOVA) by method and held-out vendor, shown on logarithmic scale. \textbf{(b)} Multivariate vendor structure (PERMANOVA $R^2$, permutational multivariate analysis of variance) with significance indicators (asterisks denote *$p<0.05$, **$p<0.01$, ***$p<0.001$; ns = not significant). \textbf{(c)} Age-related biological signal (semi-partial $R^2$) recovery. \textbf{(d)} Subject-level preservation (Pearson $r_{\text{pre,post}}$) on unseen vendors.}
\label{fig:lovo}
\end{figure}

\subsection{Structural Constraints}

Three methodological observations converge on a single structural limitation of the dataset: only 3 vendors with 44 perfectly nested sites. The LME random-effects covariance matrix was singular in 75\% of cases because 3 vendor levels provide insufficient information for stable REML variance-component estimation ($k\geq5$ recommended; Table~3). RELIEF's ICA found no vendor-associated components because ComBat-joint preprocessing had already removed vendor mean/scale effects, leaving insufficient residual vendor structure for ICA to detect. These constraints are inherent to the retrospective design; they do not invalidate the internal benchmark findings but bound their generalisability to similarly structured consortia.

\begin{table}[H]
\centring
\caption{LME convergence status by biomarker and cohort. FE = fixed-effects fallback (singular random-effects covariance matrix); LME = converged linear mixed-effects model. Only MD and AD achieved stable LME fits in both cohorts (4/16 cases, 25\%). All 12 FE-fallback cases produced non-singular post-harmonisation $\eta^2$ reductions ($>99.5\%$), confirming that FE centring is sufficient for vendor-effect removal even when REML fails.}
\label{tab:lme}
\begin{tabular}{llSSS}
\toprule
Cohort & Biomarker & $N$ & $\eta^2$ (pre) & $\eta^2$ (post) & Status \\
\midrule
\multirow{8}{*}{HC} 
  & T2w-CSA   & 203 & 0.0531 & 4.71e-04 & FE \\
  & GM-CSA    & 203 & 0.1777 & 4.00e-04 & FE \\
  & MTR       & 194 & 0.1618 & 3.61e-04 & FE \\
  & MTsat     & 190 & 0.3379 & 3.99e-05 & FE \\
  & FA        & 198 & 0.2701 & 1.54e-04 & FE \\
  & MD        & 198 & 0.3771 & 1.99e-04 & \textbf{LME} \\
  & AD        & 198 & 0.6223 & 4.77e-04 & \textbf{LME} \\
  & RD        & 198 & 0.2356 & 3.54e-04 & FE \\
\midrule
\multirow{8}{*}{ALL} 
  & T2w-CSA   & 267 & 0.0424 & 2.14e-05 & FE \\
  & GM-CSA    & 267 & 0.1701 & 3.08e-05 & FE \\
  & MTR       & 254 & 0.1800 & 5.89e-04 & FE \\
  & MTsat     & 249 & 0.3605 & 4.38e-05 & FE \\
  & FA        & 260 & 0.2838 & 1.69e-04 & FE \\
  & MD        & 260 & 0.4215 & 8.92e-05 & \textbf{LME} \\
  & AD        & 260 & 0.6611 & 1.02e-04 & \textbf{LME} \\
  & RD        & 260 & 0.2625 & 2.46e-04 & FE \\
\bottomrule
\end{tabular}
\end{table}

% ===================================================================
\section{Discussion}

This study provides, to our knowledge, the first systematic benchmark of statistical harmonisation methods for multi-vendor spinal cord qMRI biomarkers. Six principal findings emerge.

First, baseline vendor effects in spinal cord qMRI are large ($\eta^2=0.28$, PERMANOVA $R^2=0.32$)---comparable to or exceeding brain imaging reports [7,8,16]---and dominate biological variation by orders of magnitude. This justifies harmonisation as a default preprocessing step for multi-centre spinal cord studies---a need echoed by recent European quality initiatives [17] and methodological quality frameworks for quantitative imaging [22].

Second, all five methods effectively eliminate vendor effects to statistical negligibility. The choice among methods therefore hinges not on whether they work, but on the downstream analytic goal. For subject-level inference (individual prediction, longitudinal change detection), LME/FE or per-metric ComBat may be preferred. For pooled-group analyses (case-control separation, multivariate classification), CovBat or RELIEF offer maximal multivariate purity---motivated by evidence that machine learning can exploit residual covariance differences undetectable by univariate ComBat [8], and that latent scanner structure may persist after mean/scale removal [9].

Third, the Pareto frontier between decontamination and preservation is reproducible across cohorts and robust to vendor imbalance (as confirmed by Design C). This provides a quantitative, actionable decision framework.

Fourth, LOVO external validation differentiates methods sharply: ComBat and ComBat-joint generalise without degradation to unseen vendors, while RELIEF and LME/FE fail. For consortia planning to add new vendor scanners, ComBat provides a robust default choice.

Fifth, the structural constraints revealed by this benchmark---vendor-site nesting, LME singular covariance in 75\% of cases, and RELIEF ICA inactivation---reflect the 3-vendor, perfectly-nested spine-generic design. These are not methodological failures but intrinsic dataset properties that bound interpretability: reported vendor effects cannot be partitioned into vendor-level versus site-level components, and LME/RELIEF's degradation in low-vendor-count settings is expected given their modelling assumptions. Larger consortia with more vendors and crossed designs would likely resolve these constraints.

Sixth, the biological-signal amplification (2.3--3.7-fold for age) demonstrates that harmonisation does not merely remove noise---it unmasks signal previously obscured by vendor variance. This may have implications for clinical studies seeking to detect disease-related biomarker changes across sites.

\textbf{Limitations.} (i) Vendor-site perfect nesting precludes decomposition of vendor-level versus site-level effects; our findings reflect the composite ``vendor-in-site'' effect relevant to practical consortia. (ii) The dataset covers only three vendors with substantial imbalance (Siemens 67\%); extrapolation to other manufacturers requires caution. (iii) LME converged only for MD and AD (25\%); the reported LME/FE preservation advantage ($r=0.84$) may reflect the minimal correction of the FE fallback rather than LME shrinkage. (iv) RELIEF's ICA inactivation ($q\geq0.99$ for all components) may be dataset-specific; Zhang et al.\ [9] demonstrated RELIEF's effectiveness in a three-site DTI study where latent vendor structure persisted after ComBat preprocessing---suggesting that consortia with stronger residual vendor-covariance structure may benefit from RELIEF, whereas our 44-site design with ComBat-joint preprocessing may have already exhausted this structure. (v) LOVO validation used z-score matching for unseen-vendor transformation---the standard ComBat-from-template approach [7]---but does not test prospective acquisition on genuinely new hardware. (vi) This analysis was not pre-registered; however, all harmonisation methods were selected and evaluated on a pre-specified endpoint framework prior to analysis, and the complete analytical pipeline is publicly documented in the accompanying code repository, ensuring full computational reproducibility.

\textbf{Conclusion.} Statistical harmonisation effectively eliminates vendor effects in multi-centre spinal cord qMRI and is recommended as a default preprocessing step. ComBat provides the best balance of vendor removal, generalisability to unseen manufacturers, and subject-level preservation, and is recommended for broad multi-centre deployment.

\textbf{Evidence level.} Moderate — This study constitutes a large retrospective observational multicentre study (44 sites, 3 vendors, 267 subjects) with controlled design, internal validation (dual-cohort and balanced-subsample analyses), and an appropriate spectrum of cases, following the Insights into Imaging evidence framework [23].

% ===================================================================
\section*{References}

\begin{enumerate}[nosep,leftmargin=*]
\item Cohen-Adad J, Alonso-Ortiz E, Abramovic M, et al. Open-access quantitative MRI data of the spinal cord and reproducibility across participants, sites and manufacturers. Sci Data. 2021;8:219. DOI: 10.1038/s41597-021-00941-8
\item Freund P, Boller V, Emmenegger TM, et al. MRI in degenerative cervical myelopathy. Lancet Neurol. 2024;23(1):93--107. DOI: 10.1016/S1474-4422(24)00113-4
\item Valosek J, Bédard S, Gros C, et al. Normative spinal cord morphometry. NeuroImage. 2024;291:120589. DOI: 10.1016/j.neuroimage.2024.120589
\item Cohen-Adad J, Alonso-Ortiz E, Abramovic M, et al. Generic acquisition protocol for quantitative MRI of the spinal cord. Nat Protoc. 2021;16:4611--4632. DOI: 10.1038/s41596-021-00588-0
\item Prados F, Ashburner J, Blaiotta C, et al. Spinal cord grey matter segmentation challenge. NeuroImage. 2017;152:312--329. DOI: 10.1016/j.neuroimage.2017.03.010
\item Johnson WE, Li C, Rabinovic A. Adjusting batch effects in microarray expression data using empirical Bayes methods. Biostatistics. 2007;8(1):118--127. DOI: 10.1093/biostatistics/kxj037
\item Fortin JP, Cullen N, Sheline YI, et al. Harmonization of cortical thickness measurements across scanners and sites. NeuroImage. 2018;167:104--120. DOI: 10.1016/j.neuroimage.2017.11.024
\item Chen AA, Beer JC, Tustison NJ, et al. Mitigating site effects in covariance for machine learning in neuroimaging data. Hum Brain Mapp. 2022;43(4):1179--1195. DOI: 10.1002/hbm.25688
\item Zhang R, Oliver LD, Voineskos AN, et al. RELIEF: A structured multivariate approach for removal of latent inter-scanner effects. Imaging Neurosci. 2023;1:1--16. DOI: 10.1162/imag_a_00011
\item Bayer JMM, Thompson PM, Ching CRK, et al. Site effects how-to and when: an overview of retrospective techniques to accommodate site effects in multi-site neuroimaging analyses. NeuroImage. 2022;251:119009. DOI: 10.1016/j.neuroimage.2022.119009
\item Middleton DM, Li Y, Chen J, et al. ComBat harmonization for spinal cord DTI. NMR Biomed. 2023;36:e5060. DOI: 10.1002/nbm.5060
\item Jodoin PM, Gros C, Cohen-Adad J. Best practices for ComBat in diffusion MRI. NeuroImage. 2025;297:120752. DOI: 10.1016/j.neuroimage.2024.120752
\item De Leener B, Lévy S, Dupont SM, et al. SCT: Spinal Cord Toolbox, an open-source software for processing spinal cord MRI data. NeuroImage. 2017;145:24--43. DOI: 10.1016/j.neuroimage.2016.10.009
\item Anderson MJ. A new method for non-parametric multivariate analysis of variance. Austral Ecol. 2001;26:32--46. DOI: 10.1111/j.1442-9993.2001.01070.pp.x
\item De Stefano N, Battaglini M, Pareto D, et al. MAGNIMS recommendations for spinal cord MRI. Lancet Neurol. 2022;21:15--17. DOI: 10.1016/S1474-4422(21)00411-8
\item Fortin JP, Parker D, Tun\c{c} B, et al. Harmonization of multi-site diffusion tensor imaging data. NeuroImage. 2017;161:149--170. DOI: 10.1016/j.neuroimage.2017.08.047
\item Busoni S, et al. The MR quality landscape in Europe. Insights Imaging. 2026;17:125. DOI: 10.1186/s13244-026-02280-x
\item Da-Ano R, Visvikis D, Hatt M. Harmonization strategies for multicenter radiomics investigations. Phys Med Biol. 2020;65:24TR02. DOI: 10.1088/1361-6560/aba798
\item Beer JC, Tustison NJ, Cook PA, et al. Longitudinal ComBat: a method for harmonizing longitudinal multi-scanner imaging data. NeuroImage. 2020;220:117129. DOI: 10.1016/j.neuroimage.2020.117129
\item Cabini RF, Brero F, Lancia A, et al. Preliminary report on harmonization of features extraction process using the ComBat tool in the multi-center \textquotedblleft Blue Sky Radiomics\textquotedblright\ study on stage III unresectable NSCLC. Insights Imaging. 2022;13:38. DOI: 10.1186/s13244-022-01171-1
\item van Timmeren JE, Cester D, Tanadini-Lang S, et al. Radiomics in medical imaging---\textquotedblleft how-to\textquotedblright\ guide and critical reflection. Insights Imaging. 2020;11(1):91. DOI: 10.1186/s13244-020-00887-2
\item Kocak B, Akinci D'Antonoli T, Mercaldo N, et al. MEthodological RadiomICs Score (METRICS): a quality scoring tool for radiomics research endorsed by EuSoMII. Insights Imaging. 2024;15:8. DOI: 10.1186/s13244-023-01572-w
\item Mart\'i-Bonmat\'i L. Evidence levels in radiology: the insights into imaging approach. Insights Imaging. 2021;12:45. DOI: 10.1186/s13244-021-00995-7
\end{enumerate}

\end{document}
